% abstract, choose between abstract and summary
% Grounded on v1.5.1 results (see chapter5.tex, chapter6.tex, docs/ablation/RESULTS.md).
% English (matches \english in thesis.tex and the body chapters).

\begingroup
\setlength{\parindent}{0pt}
\setlength{\parskip}{\baselineskip}

Modern enterprises manage data through physical database schemas whose meaning is scattered across unstructured governance documentation---data dictionaries, policy files, and design documents. This \emph{semantic gap} between documentation and schema makes data discovery and governance slow, manual, and error-prone. This thesis presents \textbf{SemanticMesh}, a multi-agent framework that automates semantic discovery by bridging that gap: it ingests governance documents together with DDL, constructs a queryable Knowledge Graph in Neo4j that links business concepts to physical tables, and answers natural-language questions over the resulting graph through retrieval-augmented generation.

SemanticMesh is organised around a CQRS-inspired two-graph architecture: a write-heavy \emph{Builder} graph that extracts triplets, resolves entities, maps schemas to an ontology, and generates and executes Cypher; and a read-heavy \emph{Query} graph that performs hybrid retrieval. The Builder combines a two-stage entity-resolution pipeline---BGE-M3 vector blocking with Union-Find clustering followed by a mid-tier LLM judge---with self-reflection loops for quality assurance: Actor-Critic schema-mapping validation, Cypher healing with a deterministic parameterised fallback, and a Self-RAG-style hallucination grader. The Query graph fuses dense vector search, BM25 full-text search, and variable-depth graph traversal through Reciprocal Rank Fusion, then re-ranks candidates with a \texttt{bge-reranker-v2-m3} cross-encoder behind a retrieval-quality gate.

The system was evaluated through a systematic ablation campaign of 21 single-variable studies plus the optimised configuration, across seven synthetic datasets (111 tables, 210 questions) judged by a domain-specific, architecture-aware LLM-as-a-Judge. The best configuration achieved \textbf{210/210 grounded answers (100\%) with zero hallucinations} and an average judge score of \textbf{4.31/5}, including a 58-table stress set. The ablation evidence informs the default configuration and revises two intuitive expectations: on the baseline dataset most single-variable changes are indistinguishable from the baseline, so components are retained for robustness on degraded schemas rather than for measurable quality; and retrieval recall is decoupled from judged answer quality, since widening the reranker pool from $k{=}5$ to $k{=}20$ raised Gold-Standard coverage from $0.860$ to $0.986$ without raising the score (4.31 vs.\ 4.28). The thesis documents these findings, the system's observed and design-extrapolated limitations, and directions for future work centred on community detection for global reasoning and domain-specific embedding.

\endgroup
